\chapter{Accessibility, internationalization, ethics, sustainability, and AI assistance}

\section{People are part of the operating environment}
Software affects people with different abilities, languages, devices, resources,
and power. Human requirements are not decorative polish. A service that cannot
be used with a keyboard, that assumes one date format, or that quietly denies a
group access has a functional defect for those users.

WCAG 2.2 organizes accessibility guidance around perceivable, operable,
understandable, and robust principles, with testable success criteria
\cite{w3c2023wcag}. Treat conformance as a useful baseline, not as proof that a
particular person can complete a task. Test with assistive technologies and with
people who use them when possible.

\section{Practical accessibility}
For a web interface, check semantic structure, heading order, labels, focus
visibility, keyboard operation, contrast, motion preferences, error recovery,
live-region announcements, and responsive zoom. Do not rely on color alone. Make
the primary task possible without a mouse and without a fast connection.

Accessibility must include operational interfaces. A dashboard with a tiny red-
green status indicator can exclude operators with color-vision differences and
can encourage unsafe decisions under pressure.

\section{Internationalization}
Separate translatable content from code, support plural rules, keep user locale
and time zone explicit, format numbers and dates according to context, and avoid
assuming names, addresses, or writing direction. Store canonical instants and
localize presentation. Test long strings, right-to-left layouts, non-ASCII input,
and culturally unfamiliar calendars where relevant.

\section{Sustainability}
Software consumes energy through computation, storage, networks, manufacturing,
and replacement cycles. A smaller carbon footprint can come from efficient
algorithms, right-sized infrastructure, caching, data retention limits, longer
hardware life, and avoiding work that has no user value. Optimization may trade
energy against latency or availability; measure the relevant system boundary
instead of claiming that one technique is always greener.

\section{Ethics and power}
Ask who benefits, who bears risk, who can contest a decision, and who can leave
the system. A model may be statistically accurate and still unfairly allocate
opportunity. A monitoring feature may improve security and also enable coercion.
Document purpose, limits, affected groups, appeal paths, and governance. Ethical
analysis is not solved by adding the word “responsible” to a project name.

\section{Responsible AI-assisted development}
AI assistants can accelerate drafting, search, and exploration. They can also
invent APIs, reproduce insecure code, leak confidential context, or create an
illusion of review. Treat generated code as an untrusted contribution:

\begin{enumerate}
  \item specify the behaviour and constraints before prompting;
  \item review generated code like an external patch;
  \item run tests, static analysis, dependency and secret checks;
  \item verify licenses, provenance, privacy, and security-sensitive logic;
  \item retain human ownership for design, approval, and incident response.
\end{enumerate}

Never accept a citation, benchmark, command, or library interface merely because
it sounds plausible. Check primary documentation and execute safe examples.
AI assistance is a tool for increasing feedback speed, not a transfer of
accountability.

\section{Inclusive design review}
Include affected perspectives before implementation, not only in a final audit.
Use a short risk register with severity, evidence, mitigation, owner, and review
date. Ask what happens to a person who has no high-end device, cannot hear an
alert, has a different name format, or is wrongly classified. These questions
often reveal ordinary usability and reliability improvements as well as ethical
risks.

\begin{exercise}
Choose a feature flag that automatically limits accounts after suspicious
activity. Identify accessibility, language, privacy, fairness, and appeal risks.
Design a human-review path and three metrics that could reveal harm.
\end{exercise}

\section{Chapter summary}
Design for real variation in people and contexts. Use accessibility standards as
a baseline, make locale assumptions explicit, measure sustainability trade-offs,
analyze power and contestability, and verify AI-assisted work with the same or
stronger discipline as human-authored code.

